\documentclass{article}
\usepackage{amsmath, amssymb, amsthm}
\usepackage{hyperref}
\usepackage{geometry}
\geometry{margin=1in}

\title{Erd\H{o}s Problem \#912}
\date{Last updated: 2025-08-31}
\author{Source: erdosproblems.com}

\begin{document}
\maketitle

\noindent\textbf{Status:} OPEN \quad \textbf{No prize}

\noindent\textbf{Tags:} number theory, factorials

\medskip

\noindent\textbf{Problem Statement:}

If\[n! = \prod_i p_i^{k_i}\]is the factorisation into distinct primes then let $h(n)$ count the number of distinct exponents $k_i$. Prove that there exists some $c&#62;0$ such that\[h(n) \sim c \left(\frac{n}{\log n}\right)^{1/2}\]as $n\to \infty$.




A problem of Erd\H{o}s and Selfridge, who proved (see \cite{Er82c})\[h(n) \asymp \left(\frac{n}{\log n}\right)^{1/2}.\]A heuristic of Tao using the Cram\'{e}r model for the primes (detailed in the comments) suggests this is true with\[c=\sqrt{2\pi}=2.506\cdots.\]




References

[Er82c] Erd\H{o}s, P., Miscellaneous problems in number theory . Congr. Numer. (1982), 25-45.

\noindent\textbf{Related OEIS sequences:} A071626


\bigskip
\noindent\small{Source: \url{https://www.erdosproblems.com/912}}
\end{document}
